% Full answers PDF: MCQ answers + worked explanations per option, then official MS clips.
% Compiler: XeLaTeX.
\documentclass[12pt,a4paper]{article}
\usepackage[margin=18mm,headheight=14pt]{geometry}
\usepackage{fontspec}
\usepackage[version=4]{mhchem}
\usepackage{graphicx}
\usepackage{xcolor}
\usepackage{booktabs}
\usepackage{fancyhdr}
\usepackage{enumitem}
\usepackage{needspace}

\setmainfont{Helvetica Neue}

\pagestyle{fancy}
\fancyhf{}
\fancyhead[L]{\small 3.5 Shapes of molecules --- answers}
\fancyhead[R]{\small CIE 9701}
\fancyfoot[C]{\thepage}
\renewcommand{\headrulewidth}{0.3pt}

\setlength{\parindent}{0pt}
\setlength{\parskip}{0.35em}

\newcommand{\msblock}[2]{%
  \par\vspace{0.8em}%
  \noindent\begin{minipage}{\textwidth}%
  {\normalsize\bfseries #1}\par\vspace{0.25em}%
  \noindent\includegraphics[width=\textwidth,keepaspectratio]{#2}\par
  \end{minipage}\par
}

\newcommand{\mcq}[4]{%
  \par\needspace{5\baselineskip}\vspace{0.7em}%
  {\normalsize\bfseries #1.\quad #2}%
  \hfill{\footnotesize\color{gray}#3}\par\vspace{0.15em}%
  {\small #4}\par
}

\newcommand{\why}[2]{\par\needspace{2.2\baselineskip}\vspace{0.1em}{\small\textbf{#1} #2}\par}

\begin{document}

\begin{center}
{\large\bfseries 3.5 Shapes of molecules homework --- answers}\\[0.2em]
{\small CIE 9701 AS Chemistry \quad·\quad 8 multiple-choice + 10 written marks}
\end{center}

{\small\bfseries Do not issue this file to students.}
\hrule
\vspace{0.4em}

\section*{Section A --- Multiple choice}

{\normalsize\bfseries Quick answers}

\vspace{0.3em}
{\small
\begin{tabular}{@{}ll@{\hspace{2.6em}}ll@{\hspace{2.6em}}ll@{}}
\toprule
Q & Ans. & Q & Ans. & Q & Ans.\\
\midrule
A1 & \textbf{A} & A4 & \textbf{B} & A7 & \textbf{B}\\
A2 & \textbf{D} & A5 & \textbf{B} & A8 & \textbf{D}\\
A3 & \textbf{C} & A6 & \textbf{C} & & \\
\bottomrule
\end{tabular}}

\vspace{0.6em}
{\normalsize\bfseries Worked explanations}

\mcq{A1}{A}{2025 MJ Paper 12 Q7}{%
The bond angles to use are \ce{H2O} 104.5\textdegree, \ce{NH3}
107\textdegree\ and \ce{BF3} 120\textdegree. So the \emph{smallest} angle belongs
to \ce{H2O} (two lone pairs on O squeeze the H--O--H angle most) and the
\emph{largest} to \ce{BF3} (three bond pairs, no lone pair, 120\textdegree). Row
A gives \ce{H2O} as smallest and \ce{BF3} as largest.}

\why{B / C / D:}{each puts the species in the wrong order --- e.g.\ \ce{H2O} is
never the largest angle, because its two lone pairs reduce the angle below the
tetrahedral value.}

\mcq{A2}{D}{2023 MJ Paper 13 Q5}{%
\ce{NH3 + H+ -> NH4+}. In \ce{NH3}, nitrogen has 3 bond pairs \emph{and 1 lone
pair} (pyramidal, angle about 107\textdegree). When the lone pair is used to bond
to \ce{H+}, the nitrogen in \ce{NH4+} has \textbf{four bond pairs and no lone
pair}. Four identical bond pairs give a \emph{tetrahedral} shape with bond angle
\textbf{109.5\textdegree}. So the row ``tetrahedral, 109.5'' is correct.}

\why{A / B:}{say the shape is pyramidal. \ce{NH4+} has no lone pair, so it is not
pyramidal.}

\why{C:}{gives the right shape but keeps the \ce{NH3} angle of
107\textdegree; removing the lone pair opens the angle to 109.5\textdegree.}

\mcq{A3}{C}{2024 MJ Paper 11 Q24}{%
The ammonium ion, \ce{NH4+}, has four bonding pairs of electrons around nitrogen
and no lone pair, so it is tetrahedral with a bond angle of
\textbf{109.5\textdegree}.}

\why{A:}{90\textdegree\ is an octahedral angle, not tetrahedral.}

\why{B:}{107\textdegree\ is the angle in \ce{NH3} (one lone pair), not in
\ce{NH4+}.}

\why{D:}{120\textdegree\ is a trigonal planar angle (e.g.\ \ce{BF3}).}

\mcq{A4}{B}{2025 MJ Paper 14 Q11}{%
\ce{PF5} has five bonding pairs around phosphorus and \emph{no lone pair}, so it
is trigonal bipyramidal. This makes statement B correct.}

\why{A:}{not all F--P--F angles are 90\textdegree. The two axial F atoms make
90\textdegree\ with the equatorial F atoms, but the equatorial F--P--F angles are
120\textdegree.}

\why{C:}{\ce{PF5} is symmetrical, so the bond dipoles cancel and the molecule has
\emph{no} overall dipole moment.}

\why{D:}{the shape is trigonal bipyramidal, not octahedral (octahedral is 6
bond pairs, e.g.\ \ce{SF6}).}

\mcq{A5}{B}{2021 MJ Paper 11 Q34}{%
(1) \ce{C2H4} (ethene) --- the H--C--H angle is about 117\textdegree\ and the
H--C=C angle about 121\textdegree, so \ce{C2H4} has angles close to
120\textdegree. \textbf{Correct.}
(2) \ce{PF5} --- trigonal bipyramidal, so the equatorial F--P--F angles are
exactly 120\textdegree. \textbf{Correct.}
(3) \ce{NCl3} --- nitrogen has 3 bond pairs and 1 lone pair, giving a pyramidal
shape with angles of about 107\textdegree; it has no 120\textdegree\ angle.
\textbf{Wrong.}
So statements 1 and 2 only are correct.}

\why{A:}{includes \ce{NCl3}, which has no 120\textdegree\ angle.}

\why{C:}{drops \ce{C2H4}, which does have a 120\textdegree\ angle.}

\why{D:}{keeps only statement 1 and loses the correct statement about
\ce{PF5}.}

\mcq{A6}{C}{2025 MJ Paper 14 Q3}{%
Compare the smallest bond angle in each pair:
A \ce{CH4} 109.5\textdegree\ vs \ce{SF6} 90\textdegree\ (smallest) --- 109.5 is
\emph{not} smaller than 90.
B \ce{CO2} 180\textdegree\ vs \ce{BF3} 120\textdegree --- 180 is not smaller
than 120.
C \ce{H2O} 104.5\textdegree\ vs \ce{H3O+} 107\textdegree\ --- 104.5 \emph{is}
smaller than 107. \textbf{Correct.}
D \ce{NH4+} 109.5\textdegree\ vs \ce{NH3} 107\textdegree --- 109.5 is not smaller
than 107.}

\why{A / B / D:}{in each case the first species has the larger (or equal) bond
angle, so the required ``smaller in the first species'' condition fails.}

\mcq{A7}{B}{2022 MJ Paper 13 Q6}{%
In hydrazine, \ce{H2N-NH2}, each nitrogen atom has \textbf{3 bond pairs and 1 lone
pair} (2 N--H bonds, 1 N--N bond, 1 lone pair). This is the same electron-pair
arrangement as in \ce{NH3}, so the angle X is the pyramidal angle of about
\textbf{107\textdegree}.}

\why{A:}{90\textdegree\ would require octahedral electron-pair geometry.}

\why{C:}{109.5\textdegree\ is the tetrahedral angle, which needs \emph{four} bond
pairs and no lone pair --- not the case at nitrogen here.}

\why{D:}{120\textdegree\ is trigonal planar (three bond pairs, no lone pair).}

\mcq{A8}{D}{2024 MJ Paper 13 Q6}{%
Borazole, \ce{N3B3H6}, is a six-membered ring of alternating N and B atoms ---
the ``inorganic benzene''. Like benzene it is \textbf{planar}, and its
delocalised $\pi$ system contains \textbf{6} $\pi$ electrons (the aromatic
6-$\pi$-electron count).}

\why{A / B:}{say borazole is non-planar; it is planar, like benzene.}

\why{C:}{gives 3 $\pi$ electrons; borazole is aromatic with 6 $\pi$ electrons.}

\newpage

\section*{Section B --- Mark scheme clips}

\msblock{B1(a)(iii)  [1]}{cie-9701-2021-on-p21-q1a-iii-ms.png}

\msblock{B2(b)(i)  [2]}{cie-9701-2024-mj-p21-q2b-i-ms.png}

\msblock{B2(b)(ii)  [2]}{cie-9701-2024-mj-p21-q2b-ii-ms.png}

\msblock{B4(e)  [2]}{cie-9701-2023-mj-p21-q1e-ms.png}

\msblock{B5(b)(ii)  [1]}{cie-9701-2025-fm-p22-q1b-ii-ms.png}

\msblock{B6(b)(i)  [2]}{cie-9701-2025-mj-p23-q4b-i-ms.png}

\end{document}
